% chapters/05/evaluasi.tex
\bagian{Asesmen dan Instrumen Penilaian}
\label{sec:asesmen-gerak}

\begin{asesmenbox}
\textbf{Kedudukan asesmen bab ini.} Bab ini menuntaskan \SUBCPMK{1.2},
sehingga ketiga indikatornya diperiksa melalui dua instrumen: penilaian
artefak program (Tugas Praktikum 03) dan soal konsep.

\vspace{0.4em}
Materi Bab~V termasuk bahan \textbf{Kuis 1} pada minggu kelima (Materi 1--4)
dan bahan \textbf{Ujian Tengah Semester} pada minggu kesepuluh.
\end{asesmenbox}

\subsection*{Instrumen 1: Rubrik Praktikum 03}

\begin{table}[htbp]
\centering
\caption{Rubrik penilaian Praktikum 03 beserta indikator yang diukur}
\label{tab:rubrik-praktikum-03}
{\setlength{\tabcolsep}{4pt}\footnotesize%
\begin{tabularx}{\textwidth}{@{}L{3.2cm} C{1.3cm} C{2.1cm} Y@{}}
\toprule
\headrow \textbf{Kriteria} & \textbf{Bobot} & \textbf{Indikator} & \textbf{Deskripsi capaian} \\
\midrule
Gerak berbasis vektor dan selang waktu & 35\% & \ind{a} &
Objek digerakkan dengan \code{Vector2}, perpindahan dihitung dari kecepatan
dikalikan selang waktu, dan vektor arah dinormalkan sehingga gerak diagonal
tidak lebih cepat. \\
\addlinespace
Konsistensi kecepatan & 35\% & \ind{b} &
Kecepatan terukur sama pada laju bingkai 30 dan 60; tabel pengukuran pada
lembar kerja terisi dan disertai kesimpulan yang tepat. \\
\addlinespace
Penyelarasan posisi dan \code{Rect} & 30\% & \ind{c} &
Posisi disimpan sebagai pecahan, \code{Rect} diselaraskan dengan pembulatan,
dan dijelaskan mengapa pemotongan bilangan bulat merugikan. \\
\bottomrule
\end{tabularx}}
\end{table}

\subsection*{Instrumen 2: Soal Konsep}

Jawablah dengan singkat. Kunci jawaban beserta indikator yang diukur
disajikan pada Tabel~\ref{tab:kunci-konsep-ch05}.

\begin{enumerate}
  \item Mengapa vektor arah perlu dinormalkan sebelum dikalikan kecepatan?
        Sertakan perhitungan singkat untuk mendukung jawaban Anda.
  \item Tuliskan baris perhitungan perpindahan objek, lalu jelaskan satuan
        setiap unsurnya.
  \item Mengapa posisi objek disimpan sebagai bilangan pecahan, sedangkan
        \code{Rect} bertipe bilangan bulat?
  \item Apa akibatnya apabila posisi objek disimpan sebagai bilangan bulat pada
        kecepatan 20 piksel per detik?
  \item Apa perbedaan \code{normalize()} dan \code{normalize_ip()}?
\end{enumerate}

\begin{table}[htbp]
\centering
\caption{Kunci soal konsep Bab V beserta indikator yang diukur}
\label{tab:kunci-konsep-ch05}
{\setlength{\tabcolsep}{4pt}\footnotesize%
\begin{tabularx}{\textwidth}{@{}C{0.8cm} C{2.1cm} Y@{}}
\toprule
\headrow \textbf{No} & \textbf{Indikator} & \textbf{Pokok jawaban} \\
\midrule
1 & \ind{a} &
Karena vektor arah dari dua tombol yang ditekan bersamaan bernilai (1, 1)
dengan panjang $\sqrt{2} \approx 1{,}414$, sehingga kecepatannya menjadi
$200 \times 1{,}414 = 283$ piksel per detik, yaitu 41 persen lebih cepat
daripada gerak lurus. Penormalan membuat panjangnya satu untuk segala arah. \\
\addlinespace
2 & \ind{a} &
\code{position = position + velocity * delta_seconds}. Kecepatan dalam piksel
per detik, selang waktu dalam detik, hasilnya perpindahan dalam piksel. \\
\addlinespace
3 & \ind{c} &
Posisi harus tepat dan berkelanjutan sehingga memerlukan pecahan agar
perpindahan kecil terakumulasi; \code{Rect} dipakai untuk menggambar dan
memeriksa tabrakan, dan letak piksel pada layar memang berupa bilangan
bulat. \\
\addlinespace
4 & \ind{c} &
Perpindahan setiap bingkai 0,333 piksel terpotong menjadi nol pada setiap
bingkai, sehingga nilainya tidak pernah terakumulasi dan objek tidak bergerak
sama sekali. \\
\addlinespace
5 & \ind{a} &
\code{normalize()} mengembalikan vektor baru dengan panjang satu dan tidak
mengubah vektor asalnya; \code{normalize_ip()} mengubah vektor asalnya. \\
\bottomrule
\end{tabularx}}
\end{table}
